\chapter{Concluding Thoughts} \label{chapter:conclusions}

\section{Broader Lessons and Limitations} \label{section:broader_lessons_and_limitations}

MLPs appear reliable for ranking interfacial candidates when the bonding environments and geometries fall within the training domain. Screening workflows benefit most from MLPs in systems involving Ge, Sn, and well-registered C|X structures, where rank preservation is consistently strong. In contrast, strained or covalently dominated systems like Si|Ge or Si|C exhibit erratic performance, especially when interfacial bonding is asymmetric or poorly coordinated.

These results highlight the need for delta-correction strategies or domain-specific retraining, particularly for silicon-rich or mixed-interface systems. Hybrid workflows that combine MLP-based pre-screening with selective DFT confirmation offer a scalable and pragmatic approach to large-scale interface discovery.

This study reinforces the value of rank-based benchmarking — such as Spearman, Kendall, and Top-N overlap — over absolute energy metrics. In screening contexts, the ability to reliably identify top-performing configurations is more actionable than minimising raw energy error.

\section{Concluding Remarks} \label{section:concluding_remarks}

This work demonstrates that while general-purpose MLPs like MACE can successfully preserve interface rankings in selected systems, their predictive utility is highly system- and geometry-dependent. Strong performance was observed in Ge, Sn, and C|X systems with good atomic registry, as well as in chemically uniform Si|Ge alloy configurations. In contrast, Si-rich and poorly coordinated interfaces exhibited systematic rank failures, underscoring the limits of current models in strained or covalently dominated regimes.

By enabling fast, large-scale screening of interfacial structures in favourable systems, this study lays the foundation for accelerated discovery workflows. However, widespread application across diverse material combinations will require targeted retraining, delta-learning, or integrated MLP-DFT hybrid schemes. Future work should continue to refine these approaches, using rank-preservation metrics as a primary benchmark for practical utility in materials design.

Despite the promising results in rank preservation for certain systems, this study has also highlighted substantial challenges in building a truly generalisable machine-learned potential for interface prediction. Limitations such as incomplete training coverage, extrapolation to novel configurations, and model underperformance in strained or covalently bonded systems underscore the need for further refinement. These risks are being actively addressed through targeted retraining, delta-learning corrections, and expanded benchmark testing, as outlined in Chapter~\ref{chapter:next_steps}. While strong performance was observed in Ge, Sn, and alloy systems, achieving consistent accuracy across a broader range of interfaces will require careful design of training data and model architecture.